% The guide's index of supplemental files (the primary PDF first), with the titles and legends of the items outside it
% (Gate C). Everything inside "uncounted" is a supplemental item legend: the Cell word count excludes it.
\noindent\textbf{Document S1.} Figures S1--S7 and Table S1.\par
\begin{uncounted}
\medskip
\noindent\textbf{Table S2. Every declared test the paper reports, related to Figures 1--7.} Each test with its family, where the
paper reports it, its statistic and n, the estimate, \Pv{}, \Pv{} adjusted within the complete declared family, the family's
size and correction, and whether it is primary or post hoc (Excel).\par
\smallskip
\noindent\textbf{Table S3. The run registry, related to Figures 1--7.} Every training and evaluation run with its configuration,
its record and its outcome (Excel).\par
\smallskip
\noindent\textbf{Table S4. The single-cell panels, related to Figure 4.} The channel and marker genes with their FlyBase
identifiers, and their values in every atlas unit (Excel).\par
\smallskip
\noindent\textbf{Data S1. Source data for Figures 1--7 and S1--S7.} The values behind every panel, one table per figure
(.csv).\par
\smallskip
\noindent\textbf{Video S1. The trained connectome plans, move by move, related to Figure 1.} The final brain and readout solve
the hardest planning map in its \val{L18} certified moves. Left, the board and the readout's choice among the four moves (the
softmax of its logits; the chosen move white). Right, the unfolded CNS, with every cell that spikes in the move's \val{G12}-ms
run in its group color. Below, live charts of the cells spiking per move, stacked by group, and of the readout's confidence in
its move, with box pushes marked.\par
\smallskip
\noindent\textbf{Video S2. Write, hold, read: one memory trial, millisecond by millisecond, related to Figure 2.} Scene 1,
readout trained, \val{G06}-ms delay with four input windows. Left and center, the cells spiking after cue 0 and after cue 1;
right, the trace (adaptation, cue 1 minus cue 0) with the cells spiking differently in white; small boards show what the model
sees. Live charts show, per ms, the cells spiking differently and the cells holding the trace, and the trace's strength against
the adaptation's own decay (dashed), which it follows exactly once no cell spikes differently.\par
\smallskip
\noindent\textbf{Video S3. One trace, two readings, related to Figure 5.} The trained \val{G06}-ms and the untrained
\val{H12}-ms delay side by side, phase by phase, the longer delay played faster. Each panel shows the trace with the cells
spiking differently. Live charts show the cosine between the two trials' traces at matching moments, each trace's strength,
and each cue's push toward its target at the first decision. The end card shows the decision-time patterns after \val{G06},
\val{G01} and \val{H12}~ms with their cosines to the trained pattern.\par
\smallskip
\noindent\textbf{Video S4. Lost access, lost repair, related to Figure 7.} Scene 1, readout trained, lesion dose 1. First the
regimes of \fig{7B} (healthy, the lesion after learning, learning under the lesion, relearning under it and late
restoration), then the full state moved between brain states at the end of the delay, as in \fig{7J} (white frame). Each
column shows the cue-1 history's spikes, the brain state at learning and test or at writing and reading, and both cues'
histories on the board, ending in a solved or failed mark. Live charts show the cells spiking per ms and, per cue, the
readout's probability of a correct move at each decision.\par
\end{uncounted}
